\documentclass[10pt,a4paper]{amsart}
\usepackage[margin=19mm]{geometry}
\usepackage{amsmath,amssymb,mathtools}
\usepackage[scr=rsfs,cal=euler]{mathalfa}
\usepackage{xfrac}
\usepackage[unicode,hidelinks,bookmarks=false]{hyperref}
\newcommand{\ie}{i.e.\ }
\newcommand{\cf}{cf.\ }
\newcommand{\resp}{respectively\ }
\newcommand{\Subsection}{\subsection}
\providecommand{\nobreakdash}{\nobreak\hbox{-}}
\setlength{\emergencystretch}{2em}
\multlinegap0pt
\newcommand{\KBEeq}{\frac{d\phi_t}{dt}=\tilde{Q}_{t}^{\leftarrow} \phi_t, \quad \phi_0(\bfx)=\psi(\bfx), \quad \text{for } t \in [0, T].}
\usepackage{amssymb}
\usepackage{mathtools}
\usepackage{xcolor}
\usepackage{cancel}
\usepackage{tcolorbox}
\usepackage{bm}
\newtheorem{lemma}{Lemma}
\newcommand{\E}{\mathbb{E}}
\newcommand{\bfY}{{\bf Y}}
\newcommand{\bfx}{{\bf x}}
\newcommand{\bfy}{ {\bf y}}
\def\eqref#1{equation~\ref{#1}}
\title{Vocabulary-size-independent Convergence of Discrete Diffusion Models: adjoint equations induce the right space \vspace{1em}}
\author{Kelvin Kan, Xingjian Li, Benjamin J. Zhang, Tuhin Sahai, Stanley Osher, Markos A. Katsoulakis}
\date{Source: arXiv:2605.17232v5 (CC BY 4.0)}
\begin{document}
\maketitle

\noindent\emph{Source and licence.} Vocabulary-size-independent Convergence of Discrete Diffusion Models: adjoint equations induce the right space \vspace{1em}. Version arXiv:2605.17232v5 (CC BY 4.0), distributed by arXiv under the Creative Commons Attribution 4.0 International licence, http://creativecommons.org/licenses/by/4.0/, as recorded in the arXiv licence metadata for this exact version and shown on the abstract page https://arxiv.org/abs/2605.17232v5. Full licence notice: \texttt{licenses/arxiv-2605.17232-CC-BY-4.0.txt}.\par\smallskip

\noindent\textit{Editorial scope.} Counted unit: the article's lemma lemma:phi\_bound (source0.tex lines 1778--1783). The article's complete original proof of it is delivered with it (source0.tex lines 1785--1839, one \texttt{proof} environment of 55 lines). No mathematics is written, repaired or re-derived: the statement and the proof are the article's own lines, in the article's own notation, and no retained line is substituted or re-flowed.  Retained uncounted as the source-local context the proof needs: lines 904--906.  Nothing was omitted from any retained span, so the package carries no omission record.  No retained line contains a citation, so the wrapper carries no bibliography and no bibliography entry.  No retained line contains a figure environment, an image-inclusion command or drawing code, so no image is delivered and the package declares no asset dependency.  Editorial formatting only, in addition: an AMS \texttt{amsart} preamble that restates the article's own declarations and macros used by the retained lines, namely the environments \texttt{lemma} on separate counters, together with the standard packages those lines need.  The retained environments print in this document as Lemma 1 in order of appearance, whereas the article prints the same items with the numbering of its own full document.  The complete native source of the article as distributed in the arXiv e-print for this exact version is bundled unchanged as \texttt{sources/200781/source0.tex}.
\begin{equation}\label{eq:KBE}
    \KBEeq
\end{equation}

\begin{lemma}[Sup-norm Non-expansiveness Under CTMC Generators]\label{lemma:phi_bound}
    Suppose $\| \psi \|_\infty < \infty$, then the solution to the Kolmogorov backward equation~\eqref{eq:KBE} $\phi_t$ satisfies
    \begin{equation}
        \| \phi_t \|_\infty \leq \| \psi \|_\infty \quad \text{for all } t\in [0, T].
    \end{equation}
\end{lemma}

\begin{proof}
We provide two proofs here for completeness. The first direction is using a Feynman-Kac representation, which is short and clean. The second direction is to mimic the maximum principle used for parabolic partial differential equations (PDEs), which is our original intuition in deriving this result.

\textbf{First Direction (Feynman-Kac representation):}
The Feynman--Kac representation gives $\phi_t(\bfx) = \E^\bfx[\psi(\bfY_t)]$, where $(\bfY_s)$ is the CTMC with generator $\tilde{Q}_{t}^{\leftarrow}$ started at $\bfY_0 = \bfx$. Hence $|\phi_t(\bfx)| \le \E^\bfx[|\psi(\bfY_t)|] \le \|\psi\|_\infty$.

\textbf{Second Direction (Maximum principle):} Recall that the Kolmogorov backward equation is given by
    \begin{equation*}
    \frac{d\phi_t}{dt}=\tilde{Q}_{t}^{\leftarrow} \phi_t, \quad \phi_0(\bfx)=\psi(\bfx), \quad \text{for } t \in [0, T].
    \end{equation*}
    For any $\epsilon >0$, construct the function
    \begin{equation*}
        \phi_t^\epsilon(\bfx) := \phi_t(\bfx) - \epsilon t; \quad \text{in vector notation: } \phi_t^\epsilon := \phi_t - \epsilon t\mathbf{1}.
    \end{equation*}
    Evaluating its time derivative gives
    \begin{align}
        \frac{d \phi_t^\epsilon }{dt} &= \frac{d}{dt}(\phi_t - \epsilon t \mathbf{1}) \nonumber \\
        &= \tilde{Q}_{t}^{\leftarrow} \phi_t - \epsilon \mathbf{1} \nonumber \\
        &= \tilde{Q}_{t}^{\leftarrow} \left( \phi^\epsilon_t + \epsilon t \mathbf{1} \right) - \epsilon \mathbf{1} \nonumber \\
        &= \tilde{Q}_{t}^{\leftarrow}  \phi^\epsilon_t  - \epsilon \mathbf{1} ,  \label{eq:d_phi_epsilon} \end{align}
        as $\tilde{Q}_{t}^{\leftarrow}$'s rows sum to 0.
        
        Suppose that $\phi^\epsilon_t$ attains its maximum at $t_0 \in (0, T]$ and $\bfx_0$, then we have the following two properties:
        \begin{enumerate}
            \item \begin{equation}\label{eq:testfunc_deri_nonneg}\frac{d  }{dt} \phi_{t_0}^\epsilon(\bfx_0) \geq 0,\end{equation} \\
            as $\frac{d  }{dt} \phi_{t_0}^\epsilon(\bfx_0)  = 0$ if $t_0 < T$ and $\frac{d  }{dt} \phi_{t_0}^\epsilon(\bfx_0)  \geq 0$ if $t_0 = T$,
            \item \begin{equation}\label{eq:testfunc_Q_neg}[\tilde{Q}_{t_0}^{\leftarrow} \phi_{t_0}^\epsilon](\bfx_0) \leq 0,\end{equation} \\
            because
            \begin{align*}
                [\tilde{Q}_{t_0}^{\leftarrow} \phi_{t_0}^\epsilon](\bfx_0) &= \sum_\bfy \tilde{Q}_{t_0}^{\leftarrow}(\bfx_0, \bfy) \phi_{t_0}^\epsilon(\bfy)\\
                &=\sum_{\bfy \neq \bfx_0} \left[\tilde{Q}_{t_0}^{\leftarrow}(\bfx_0, \bfy) \phi_{t_0}^\epsilon(\bfy) \right] + \tilde{Q}_{t_0}^{\leftarrow}(\bfx_0, \bfx_0) \phi_{t_0}^\epsilon(\bfx_0)\\
                &=\sum_{\bfy \neq \bfx_0} \underbrace{\tilde{Q}_{t_0}^{\leftarrow}(\bfx_0, \bfy)}_{\geq 0} \underbrace{\left(\phi_{t_0}^\epsilon(\bfy)-\phi_{t_0}^\epsilon(\bfx_0) \right)}_{\leq 0, \text{ as $\phi_{t_0}^\epsilon(\bfx_0)$ is max}} \leq 0.
            \end{align*}
        \end{enumerate}
    The two properties together give us
    \begin{align*}
        0 &\leq \frac{d  }{dt} \phi_{t_0}^\epsilon(\bfx_0) , \quad \text{by~\eqref{eq:testfunc_deri_nonneg}}\\
        &= \underbrace{\left[ \tilde{Q}_{t}^{\leftarrow}  \phi^\epsilon_t\right](\bfx_0)}_{\leq 0, \text{ by~\eqref{eq:testfunc_Q_neg}}}  - \epsilon ,\quad \text{by~\eqref{eq:d_phi_epsilon},}\\
        &\leq -\epsilon \\
        &<0,
    \end{align*}
    leading to a contradiction. This implies $\phi^\epsilon_t$ cannot attain its maximum at $t_0 \in (0, T]$. Thus, $\phi^\epsilon_t$ must attain its maximum at $t=0$, and we have, for all $t \in [0, T], \bfx$,
    \begin{equation*}
        \phi_t(\bfx)-\epsilon t= \phi_t^\epsilon(\bfx) \leq \max_\bfy \phi_0^\epsilon(\bfy) = \max_\bfy \phi_0(\bfy)= \max_\bfy \psi(\bfy).
    \end{equation*}
    Since $\epsilon>0$ is arbitrary, taking $\epsilon \to 0^+$ gives
    \begin{equation*}
        \phi_t(\bfx) \leq \max_\bfy \psi(\bfy), \quad \text{for all $t \in [0, T], \bfx$}.
    \end{equation*}
    This implies 
    $$
    \| \phi_t \|_\infty \leq \| \psi \|_\infty, \quad \text{for all $t \in [0, T]$},
    $$
    as desired.
\end{proof}

\end{document}
